If the lines are parallel, denoting $d$ their distance,
if $d > a$, then the locus of points is empty.
If $d \leq a$, the locus of points is constituted by the two parallel lines at distance $\mfrac{1}{2} \paren{a - d}$ from the given lines, externally.

\medskip
If the two lines intersect, then the locus of points is (the boundary of) a rectangle,
whose vertices are the two points on $r$ at distance $a$ from $s$,
and the two points on $s$ at distance $a$ from $r$,
and, in particular, the diagonals of the rectangle lie on the lines $r$ and $s$.

Those four points are clearly part of the locus,
since they are at distance $0$ from one line and at distance $a$ from the other.

Now to prove that it is a rectangle:

TODO


Now to prove that the sides are in the locus:

TODO


Now to prove that every point of the locus is on one of the sides of the rectangle:
    In fact, the interior of the rectangle is the set of points such that the sum of distances is $< a$;
    and its exterior is the set of points such that the sum of distances is $> a$.

    Prove that the points in the interior of the rectangle have $< a$:

    TODO


    Prove that the points outside have $> a$:

    TODO
